%
% A Z model of the PR review state in standards/pr-review.md: the review
% markers (reviewed_head, delta_head), the Qodo and Bugbot rounds, Bugbot
% runs and rulings, findings and answered(f), escalations, and merge_gate.
%
% Type-check:  fuzz -t standards/pr-review-model.tex
% Model-check: probcli standards/pr-review-model.tex -model_check \
%                -card COMMIT 3 -card FINDING 2 -card RUN 2 \
%                -p MAX_INITIALISATIONS 1 -p SYMMETRY_MODE hash
%   (1.77M states, about 8 minutes). Without symmetry, a precise run:
%   -card COMMIT 3 -card FINDING 1 -card RUN 1 -p MAX_INITIALISATIONS 10
% probcli writes pr-review-model.fuzz beside the file; do not commit it.
%
% A review event raises at most one finding (#new? <= 1); two findings
% raised at once behave as two events in sequence.
%
% The demo gate (demo_ok), Copilot, CI, and bot-authored (skipped) PRs are
% out of scope: each is a conjunct of merge_gate that no other term reads.
%
\documentclass{article}
\usepackage{fuzz}

\begin{document}

\title{A Formal Model of the PR Review State (\texttt{pr-review.md})}
\author{Punt Labs}
\maketitle

\section{Purpose and scope}

Eight empirical fix rounds on the review-state rules each surfaced the next
defect. This model replaces those rounds: it states the state the leader reads
from GitHub, the events that change it, and the guard $merge\_gate$ puts on
$Merge$. ProB checks that no reachable state merges an unreviewed head, an
unanswered finding, or a live un-ruled Bugbot run, and that every non-merged
state can still move.

Time is abstract. ``30 minutes have passed'' is the event $Timeout$; the model
checker explores it firing at every point. A fix push is a $Push$ plus an
$AnswerFinding$; the model lets them happen in either order, a superset of the
leader's behaviour.

\section{Given sets and enumerations}

Commits, findings, and Bugbot check runs are opaque identities.

\begin{zed}
  [COMMIT, FINDING, RUN]
\end{zed}

\begin{zed}
  Flag ::= yes | no
\end{zed}

A Bugbot run is live (queued or in progress) or done. A done run either
finished its review ($doneOk$: success or neutral) or reviewed nothing
($doneFail$: failure, timeout, cancelled).

\begin{zed}
  RunStatus ::= live | doneOk | doneFail
\end{zed}

\begin{zed}
  ROUND == 1 \upto 3
\end{zed}

\section{State}

\textbf{Commits.} $anc$ is ancestry, reflexive: $a \mapsto b \in anc$ when $a$
is $b$ or an ancestor of $b$. A push extends it; a rebase starts a fresh
history, so the old commits are no longer ancestors of the head. $used$ keeps
commit identities from being reused.

\begin{schema}{Commits}
  head : COMMIT \\
  used : \power COMMIT \\
  anc : COMMIT \rel COMMIT
\where
  head \in used \\
  head \mapsto head \in anc \\
  \dom anc \subseteq used \\
  \ran anc \subseteq used
\end{schema}

\textbf{Review markers.} $reviewedHead$ and $deltaHead$ are optional commits,
sets of size 0 or 1. $deltaRun$ is the head a local delta review was started
on and has not yet completed; $deltaRange$ the commits it covers; $deltaF$ the
findings the last completed delta review raised. $seen$ and $overlap$ are
ghost variables, not read by any guard: $seen$ is every commit a completed
review actually covered, and $overlap$ records that a delta review re-covered a
commit already seen.

\begin{schema}{Review}
  reviewedHead, deltaHead, deltaRun : \power COMMIT \\
  deltaRange, seen : \power COMMIT \\
  deltaF : \power FINDING \\
  overlap : Flag \\
  round : ROUND \\
  qodoDone : \power ROUND
\where
  \# reviewedHead \leq 1 \\
  \# deltaHead \leq 1 \\
  \# deltaRun \leq 1
\end{schema}

\textbf{Bugbot.} $runs$ gives each run its status, $runRound$ the round whose
trigger started it. $stuck$ marks runs live for over 60 minutes; $proceed$
marks runs the operator ruled ``proceed''. $retried$ is the rounds whose
leader posted a ``bugbot run'' re-run comment; $timedOut$ is the rounds whose
30-minute fallback window has elapsed since the trigger or the re-run comment.

\begin{schema}{Bugbot}
  runs : RUN \pfun RunStatus \\
  runRound : RUN \pfun ROUND \\
  stuck, proceed : \power RUN \\
  retried, timedOut : \power ROUND
\where
  \dom runRound = \dom runs \\
  stuck \subseteq \dom runs \\
  proceed \subseteq stuck
\end{schema}

\textbf{Findings and escalations.} $escRun$ and $escF$ name what the one open
escalation is about.

\begin{schema}{Issues}
  findings, answered, escF : \power FINDING \\
  escRun : \power RUN \\
  escalOpen : Flag
\where
  answered \subseteq findings \\
  escF \subseteq findings \setminus answered \\
  escalOpen = no \implies escRun = \emptyset \land escF = \emptyset
\end{schema}

\begin{schema}{Gate}
  merged : Flag
\end{schema}

\begin{schema}{State}
  Commits \\
  Review \\
  Bugbot \\
  Issues \\
  Gate
\end{schema}

\textbf{The invariant.} Every reachable state must satisfy these; ProB checks
them in every state it visits. They live in $Invariant$, not in $State$: a
predicate of $State$ is part of every $\Delta State$, so it would silently
disable a violating operation instead of reporting it.
(I1) A merge never leaves an unreviewed head, an unanswered finding, or a live
un-ruled run.
(I2) Every reviewed marker was actually covered by a completed review.
(I3) No delta review re-covers a commit an earlier review already covered.
(I4) A merge on the Bugbot fallback never skips the re-run of a failed run.

\begin{schema}{Invariant}
  State
\where
  merged = yes \implies \\
  \quad~ ( head \in seen \land findings \subseteq answered \land \\
  \quad~\quad~ ( \forall r : \dom runs @ runs~r = live \implies r \in proceed ) ) \\
  reviewedHead \subseteq seen \\
  overlap = no \\
  merged = yes \implies \\
  \quad~ ( ( \exists r : \dom runs @ runRound~r = round \land runs~r = doneOk ) \lor \\
  \quad~\quad~ ( ( \exists r : \dom runs @ runRound~r = round \land runs~r = doneFail ) \\
  \quad~\quad~\quad~ \implies round \in retried ) )
\end{schema}

\section{Derived terms}

These are the terms of \texttt{pr-review.md}, stated over $State$.

$bugbot\_open$: a live run with no ``proceed'' ruling.

\begin{schema}{BugbotOpen}
  State
\where
  \exists r : \dom runs @ runs~r = live \land r \notin proceed
\end{schema}

$bugbot\_in(N)$: a run of round $N$ finished its review; or the fallback
window elapsed, no run is open, and every failed run of round $N$ has had its
re-run requested. The last clause is the fix for defect (2).

\begin{schema}{BugbotIn}
  State
\where
  ( \exists r : \dom runs @ runRound~r = round \land runs~r = doneOk ) \lor \\
  ( round \in timedOut \land \\
  \quad~ ( \forall r : \dom runs @ runs~r = live \implies r \in proceed ) \land \\
  \quad~ ( ( \exists r : \dom runs @ runRound~r = round \land runs~r = doneFail ) \\
  \quad~\quad~ \implies round \in retried ) )
\end{schema}

$\lnot pending$: Qodo and Bugbot are both in for the current round.

\begin{schema}{NotPending}
  BugbotIn
\where
  round \in qodoDone
\end{schema}

\section{Initial state}

The PR opens on some commit; round 1; nothing reviewed, no run, no finding.

\begin{schema}{Init}
  State'
\where
  used' = \{ head' \} \\
  anc' = \{ head' \mapsto head' \} \\
  reviewedHead' = \emptyset \\
  deltaHead' = \emptyset \\
  deltaRun' = \emptyset \\
  deltaRange' = \emptyset \\
  seen' = \emptyset \\
  deltaF' = \emptyset \\
  overlap' = no \\
  round' = 1 \\
  qodoDone' = \emptyset \\
  runs' = \emptyset \\
  runRound' = \emptyset \\
  stuck' = \emptyset \\
  proceed' = \emptyset \\
  retried' = \emptyset \\
  timedOut' = \emptyset \\
  findings' = \emptyset \\
  answered' = \emptyset \\
  escF' = \emptyset \\
  escRun' = \emptyset \\
  escalOpen' = no \\
  merged' = no
\end{schema}

\section{Branch operations}

$Push$ adds a commit on top of the head.

\begin{schema}{Push}
  \Delta State \\
  \Xi Review \\
  \Xi Bugbot \\
  \Xi Issues \\
  \Xi Gate \\
  c? : COMMIT
\where
  merged = no \\
  c? \notin used \\
  head' = c? \\
  used' = used \cup \{ c? \} \\
  anc' = anc \cup ( \{ a : COMMIT | a \mapsto head \in anc \} \cross \{ c? \} ) \cup \{ c? \mapsto c? \}
\end{schema}

$Rebase$ replaces the history: the new head has no old ancestor, so neither
marker is an ancestor of it.

\begin{schema}{Rebase}
  \Delta State \\
  \Xi Review \\
  \Xi Bugbot \\
  \Xi Issues \\
  \Xi Gate \\
  c? : COMMIT
\where
  merged = no \\
  c? \notin used \\
  head' = c? \\
  used' = used \cup \{ c? \} \\
  anc' = anc \cup \{ c? \mapsto c? \}
\end{schema}

\section{Qodo}

Qodo's review for the current round completes on the head. It covers the
whole PR, sets $reviewedHead$ to the head, and may raise findings.

\begin{schema}{QodoCompletes}
  \Delta State \\
  \Xi Commits \\
  \Xi Bugbot \\
  \Xi Gate \\
  new? : \power FINDING
\where
  merged = no \\
  round \notin qodoDone \\
  new? \cap findings = \emptyset \\
  \# new? \leq 1 \\
  qodoDone' = qodoDone \cup \{ round \} \\
  reviewedHead' = \{ head \} \\
  seen' = seen \cup \{ a : COMMIT | a \mapsto head \in anc \} \\
  findings' = findings \cup new? \\
  deltaHead' = deltaHead \\
  deltaRun' = deltaRun \\
  deltaRange' = deltaRange \\
  deltaF' = deltaF \\
  overlap' = overlap \\
  round' = round \\
  answered' = answered \\
  escF' = escF \\
  escRun' = escRun \\
  escalOpen' = escalOpen
\end{schema}

\section{Bugbot}

A run starts, for the current round.

\begin{schema}{BugbotStart}
  \Delta State \\
  \Xi Commits \\
  \Xi Review \\
  \Xi Issues \\
  \Xi Gate \\
  r? : RUN
\where
  merged = no \\
  r? \notin \dom runs \\
  runs' = runs \cup \{ r? \mapsto live \} \\
  runRound' = runRound \cup \{ r? \mapsto round \} \\
  stuck' = stuck \\
  proceed' = proceed \\
  retried' = retried \\
  timedOut' = timedOut
\end{schema}

A live run ends. A finished review may raise findings.

\begin{schema}{BugbotEnd}
  \Delta State \\
  \Xi Commits \\
  \Xi Review \\
  \Xi Gate \\
  r? : RUN \\
  s? : RunStatus \\
  new? : \power FINDING
\where
  merged = no \\
  r? \in \dom runs \\
  runs~r? = live \\
  s? \in \{ doneOk, doneFail \} \\
  ( s? = doneFail \implies new? = \emptyset ) \\
  new? \cap findings = \emptyset \\
  \# new? \leq 1 \\
  runs' = runs \oplus \{ r? \mapsto s? \} \\
  findings' = findings \cup new? \\
  runRound' = runRound \\
  stuck' = stuck \\
  proceed' = proceed \\
  retried' = retried \\
  timedOut' = timedOut \\
  answered' = answered \\
  escF' = escF \\
  escRun' = escRun \\
  escalOpen' = escalOpen
\end{schema}

A live run passes 60 minutes.

\begin{schema}{BugbotStuck}
  \Delta State \\
  \Xi Commits \\
  \Xi Review \\
  \Xi Issues \\
  \Xi Gate \\
  r? : RUN
\where
  merged = no \\
  r? \in \dom runs \\
  runs~r? = live \\
  r? \notin stuck \\
  stuck' = stuck \cup \{ r? \} \\
  runs' = runs \\
  runRound' = runRound \\
  proceed' = proceed \\
  retried' = retried \\
  timedOut' = timedOut
\end{schema}

The 30-minute fallback window for the current round elapses.

\begin{schema}{Timeout}
  \Delta State \\
  \Xi Commits \\
  \Xi Review \\
  \Xi Issues \\
  \Xi Gate
\where
  merged = no \\
  round \notin timedOut \\
  timedOut' = timedOut \cup \{ round \} \\
  runs' = runs \\
  runRound' = runRound \\
  stuck' = stuck \\
  proceed' = proceed \\
  retried' = retried
\end{schema}

The leader comments ``bugbot run'' once per round when a run of that round
reviewed nothing. The guard does not read $bugbot\_in$: the fallback may
already have elapsed (fix for defect (2)). The comment restarts the fallback
window.

\begin{schema}{RequestRetry}
  \Delta State \\
  \Xi Commits \\
  \Xi Review \\
  \Xi Issues \\
  \Xi Gate
\where
  merged = no \\
  ( \exists r : \dom runs @ runRound~r = round \land runs~r = doneFail ) \\
  round \notin retried \\
  retried' = retried \cup \{ round \} \\
  timedOut' = timedOut \setminus \{ round \} \\
  runs' = runs \\
  runRound' = runRound \\
  stuck' = stuck \\
  proceed' = proceed
\end{schema}

\section{Escalations and rulings}

The leader escalates a stuck run, once per run, while no escalation is open.

\begin{schema}{EscalateStuck}
  \Delta State \\
  \Xi Commits \\
  \Xi Review \\
  \Xi Bugbot \\
  \Xi Gate \\
  r? : RUN
\where
  merged = no \\
  escalOpen = no \\
  r? \in stuck \\
  runs~r? = live \\
  r? \notin proceed \\
  escalOpen' = yes \\
  escRun' = \{ r? \} \\
  escF' = \emptyset \\
  findings' = findings \\
  answered' = answered
\end{schema}

The operator rules ``proceed'': the run no longer holds the gate.

\begin{schema}{OperatorProceed}
  \Delta State \\
  \Xi Commits \\
  \Xi Review \\
  \Xi Gate
\where
  merged = no \\
  escalOpen = yes \\
  escRun \neq \emptyset \\
  proceed' = proceed \cup escRun \\
  escalOpen' = no \\
  escRun' = \emptyset \\
  escF' = escF \\
  findings' = findings \\
  answered' = answered \\
  runs' = runs \\
  runRound' = runRound \\
  stuck' = stuck \\
  retried' = retried \\
  timedOut' = timedOut
\end{schema}

The operator rules ``wait'': the run holds the gate until it ends.

\begin{schema}{OperatorWait}
  \Delta State \\
  \Xi Commits \\
  \Xi Review \\
  \Xi Bugbot \\
  \Xi Gate
\where
  merged = no \\
  escalOpen = yes \\
  escRun \neq \emptyset \\
  escalOpen' = no \\
  escRun' = \emptyset \\
  escF' = escF \\
  findings' = findings \\
  answered' = answered
\end{schema}

In round 2 the leader escalates a new High instead of answering it.

\begin{schema}{EscalateFinding}
  \Delta State \\
  \Xi Commits \\
  \Xi Review \\
  \Xi Bugbot \\
  \Xi Gate \\
  f? : FINDING
\where
  merged = no \\
  escalOpen = no \\
  round = 2 \\
  f? \in findings \setminus answered \\
  escalOpen' = yes \\
  escF' = \{ f? \} \\
  escRun' = \emptyset \\
  findings' = findings \\
  answered' = answered
\end{schema}

The operator rules on the escalated finding; the leader applies the ruling,
which answers it.

\begin{schema}{OperatorRule}
  \Delta State \\
  \Xi Commits \\
  \Xi Review \\
  \Xi Bugbot \\
  \Xi Gate
\where
  merged = no \\
  escalOpen = yes \\
  escF \neq \emptyset \\
  answered' = answered \cup escF \\
  escalOpen' = no \\
  escF' = \emptyset \\
  escRun' = \emptyset \\
  findings' = findings
\end{schema}

\section{Findings}

The leader answers a finding by its source's path: a fix, a ledger line, or a
Qodo CLI close.

\begin{schema}{AnswerFinding}
  \Delta State \\
  \Xi Commits \\
  \Xi Review \\
  \Xi Bugbot \\
  \Xi Gate \\
  f? : FINDING
\where
  merged = no \\
  f? \in findings \setminus answered \\
  f? \notin escF \\
  answered' = answered \cup \{ f? \} \\
  findings' = findings \\
  escF' = escF \\
  escRun' = escRun \\
  escalOpen' = escalOpen
\end{schema}

\section{Rounds}

Round 1 is in and a High was fixed: the leader opens round 2, the last.

\begin{schema}{OpenRound2}
  \Delta State \\
  \Xi Commits \\
  \Xi Bugbot \\
  \Xi Issues \\
  \Xi Gate
\where
  merged = no \\
  NotPending \\
  round = 1 \\
  round' = 2 \\
  reviewedHead' = reviewedHead \\
  deltaHead' = deltaHead \\
  deltaRun' = deltaRun \\
  deltaRange' = deltaRange \\
  deltaF' = deltaF \\
  seen' = seen \\
  overlap' = overlap \\
  qodoDone' = qodoDone
\end{schema}

\section{The local delta review}

The delta review starts when nothing is pending and the head is covered by
neither marker. The guard does not ask whether the head is ``newer'' than a
marker: an unset $deltaHead$ (fix for defect (1)) and a marker a rebase
orphaned (fix for defect (4)) both leave the head uncovered. The range starts
at every marker that is an ancestor of the head --- $deltaHead$ as well as
$reviewedHead$ (fix for defect (3)) --- and is the whole diff when neither is.

\begin{schema}{DeltaReviewStart}
  \Delta State \\
  \Xi Commits \\
  \Xi Bugbot \\
  \Xi Issues \\
  \Xi Gate
\where
  merged = no \\
  NotPending \\
  deltaRun = \emptyset \\
  head \notin reviewedHead \cup deltaHead \\
  deltaRun' = \{ head \} \\
  deltaRange' = \{ a : COMMIT | a \mapsto head \in anc \} \setminus \\
  \quad~ \{ a : COMMIT | ( \exists m : reviewedHead \cup deltaHead @ \\
  \quad~\quad~ m \mapsto head \in anc \land a \mapsto m \in anc ) \} \\
  ( deltaRange' \cap seen = \emptyset \implies overlap' = overlap ) \\
  ( deltaRange' \cap seen \neq \emptyset \implies overlap' = yes ) \\
  reviewedHead' = reviewedHead \\
  deltaHead' = deltaHead \\
  deltaF' = deltaF \\
  seen' = seen \\
  round' = round \\
  qodoDone' = qodoDone
\end{schema}

The delta review completes. Only now is $deltaHead$ set; its findings are
findings like any other.

\begin{schema}{DeltaReviewCompletes}
  \Delta State \\
  \Xi Commits \\
  \Xi Bugbot \\
  \Xi Gate \\
  new? : \power FINDING
\where
  merged = no \\
  deltaRun \neq \emptyset \\
  new? \cap findings = \emptyset \\
  \# new? \leq 1 \\
  deltaHead' = deltaRun \\
  deltaRun' = \emptyset \\
  seen' = seen \cup deltaRange \\
  deltaRange' = \emptyset \\
  deltaF' = new? \\
  findings' = findings \cup new? \\
  reviewedHead' = reviewedHead \\
  overlap' = overlap \\
  round' = round \\
  qodoDone' = qodoDone \\
  answered' = answered \\
  escF' = escF \\
  escRun' = escRun \\
  escalOpen' = escalOpen
\end{schema}

When a completed delta review covers the head and its findings are answered,
$reviewedHead$ moves to it. It reads $deltaHead$, which only a completed review
sets, never the review in flight (fix for defect (5)).

\begin{schema}{Promote}
  \Delta State \\
  \Xi Commits \\
  \Xi Bugbot \\
  \Xi Issues \\
  \Xi Gate
\where
  merged = no \\
  deltaHead = \{ head \} \\
  head \notin reviewedHead \\
  deltaF \subseteq answered \\
  reviewedHead' = deltaHead \\
  deltaHead' = deltaHead \\
  deltaRun' = deltaRun \\
  deltaRange' = deltaRange \\
  deltaF' = deltaF \\
  seen' = seen \\
  overlap' = overlap \\
  round' = round \\
  qodoDone' = qodoDone
\end{schema}

\section{Merge}

$merge\_gate$: nothing pending, no open run, the head reviewed, every finding
answered, at most two rounds, no open escalation.

\begin{schema}{Merge}
  \Delta State \\
  \Xi Commits \\
  \Xi Review \\
  \Xi Bugbot \\
  \Xi Issues
\where
  merged = no \\
  NotPending \\
  ( \forall r : \dom runs @ runs~r = live \implies r \in proceed ) \\
  head \in reviewedHead \\
  findings \subseteq answered \\
  round \leq 2 \\
  escalOpen = no \\
  merged' = yes
\end{schema}

After the merge nothing changes. This stutter keeps the merged state from
reading as a deadlock.

\begin{schema}{Merged}
  \Xi State
\where
  merged = yes
\end{schema}

\end{document}
